\documentclass[12pt]{article}
\usepackage[utf8]{inputenc}
\usepackage[T1]{fontenc}
\usepackage{hyperref}
\title{Northern Tarahumara Language Resources for K-12 Curriculum Development}
\author{Compiled by Claude Sonnet 5 (claude-sonnet-5) for Andr\'as Kornai}
\date{}
\begin{document}
\maketitle

\section{Introduction}
Northern Tarahumara (ISO 639-3: thh; Glottocode nort2956) is a Uto-Aztecan language of Mexico, described as dormant. A web search found no resources specific to this variety. The only documentation found concerns other Rar\'amuri varieties (Choguita Rar\'amuri, and Central Rar\'amuri from Norogachi), which must not be mistaken for Northern Tarahumara.

\section{Primary Educational Resources}
None found for Northern Tarahumara.

\section{Linguistic and Reference Materials}
These are about a different variety and are listed only as comparative background. Gabriela Caballero's \emph{A Grammar of Choguita Rar\'amuri} (Language Science Press, 2022) is based on a documentary corpus collected between 2003 and 2018:
\url{https://langsci-press.org/catalog/book/324}

\section{Religious and Community Resources}
None found in the search results.

\section{Digital and Audio Resources}
The Choguita Rar\'amuri corpus (over 200 hours) is reported to be held at the Endangered Languages Archive and at the Survey of California and Other Indian Languages. A guide to that collection is hosted at:
\url{https://scholarspace.manoa.hawaii.edu/handle/10125/24734}
It is again a different variety.

\section{Curriculum Development Context}
Northern Tarahumara is dormant, with no remaining first-language speakers according to the identification given. Any curriculum would be revitalization work and would have to start from archival and comparative material. The search found no evidence of school teaching.
\end{document}